\documentclass[10pt,a4paper]{amsart}
\usepackage[margin=19mm]{geometry}
\usepackage{amsmath,amssymb,mathtools}
\usepackage[scr=rsfs,cal=euler]{mathalfa}
\usepackage{xfrac}
\usepackage[unicode,hidelinks,bookmarks=false]{hyperref}
\newcommand{\ie}{i.e.\ }
\newcommand{\cf}{cf.\ }
\newcommand{\resp}{respectively\ }
\newcommand{\Subsection}{\subsection}
\providecommand{\nobreakdash}{\nobreak\hbox{-}}
\setlength{\emergencystretch}{2em}
\multlinegap0pt
\usepackage{bm}
\usepackage{bbm}
\usepackage{dsfont}
\usepackage{xspace}
\usepackage{cancel}
\usepackage{mathtools}
\usepackage{amsthm}
\makeatletter
\@ifundefined{thm}{\newtheorem{thm}{Theorem}}{}
\@ifundefined{example}{\newtheorem{example}[thm]{Example}}{}
\makeatother
\title[Cyclic sieving on permutations -- an analysis of maps and st]{Cyclic sieving on permutations -- an analysis of maps and statistics in the FindStat database}
\author{Ashleigh Adams, Jennifer Elder, Nadia Lafreni\`ere, Erin McNicholas, Jessica Striker, Amanda Welch}
\date{Source: arXiv:2402.16251v2 (CC BY 4.0)}
\begin{document}
\maketitle

\noindent\emph{Source and licence.} Cyclic sieving on permutations -- an analysis of maps and statistics in the FindStat database. Version arXiv:2402.16251v2 (CC BY 4.0), distributed under the Creative Commons Attribution 4.0 International licence, as recorded in the arXiv licence metadata for this exact version and shown on the abstract page https://arxiv.org/abs/2402.16251v2. Full rights notice: licenses/arxiv-2402.16251-CC-BY-4.0.txt.\par\smallskip

\noindent\textit{Editorial scope.} Counted unit: the Thm whose source label is \texttt{thm:CSP\_AK-Invisible} in the article (source0.tex lines 1208--1210, source label \texttt{thm:CSP\_AK-Invisible}), delivered together with the article's own complete original proof of it (source0.tex lines 1212--1241, one \texttt{proof} environment of 30 lines). No mathematics is written, repaired or re-derived: statement and proof are the article's own text line for line. Required context retained uncounted: thm:CSP\_AK (lines 1157--1160); ex:kappa\_psi (lines 1243--1249). Every definition, display and label of those lines is retained unchanged. Omitted from this package and not redistributed: 1--1156, 1161--1207, 1211--1211, 1242--1242, 1250--2003. Nothing inside a retained excerpt is omitted or altered: every retained body is delivered exactly as the article's lines are written, so no omission receipt of any kind accompanies this package. Citations: the retained excerpts carry no citation command at all, so no bibliography entry of the article is transcribed and no reference is redistributed. Reference adaptation: none was needed and none was made; every reference inside the delivered unit resolves inside it, so no unresolved reference or citation is produced. Printed slips kept literal: no printed word, symbol, hypothesis or number of the counted statement or of its proof was altered. Layout and class: the article's own class is \texttt{amsart} with packages including amsmath, amssymb, fullpage, amsthm, amsfonts, newlfont, url, xspace, stmaryrd, xcolor, graphics, graphicx, verbatim, float; the wrapper is the lane's pinned \texttt{amsart} editorial wrapper at 10pt on A4, which restates the theorem-like environments the retained text needs and adds the article's own macro definitions that the retained text needs (none), with extra wrapper packages (bm, bbm, dsfont, xspace, cancel, mathtools). The delivered numbering is the unit's own. Assets: no retained span contains a figure environment, a graphics-inclusion command, a PDF-page inclusion, a table or a TikZ picture; no image file is copied into this package, the package declares no asset dependency and no image file is redistributed with it. Licence: version arXiv:2402.16251v2 of this article is distributed under the Creative Commons Attribution 4.0 International licence, as recorded in the arXiv licence metadata for this exact version and shown on the abstract page https://arxiv.org/abs/2402.16251v2. A full rights notice is delivered at licenses/arxiv-2402.16251-CC-BY-4.0.txt. Characters: every retained excerpt is pure ASCII, so the delivered engine, XeLaTeX with its default Latin Modern text font, reports no missing character.
\begin{thm}\label{thm:CSP_AK}
  The statistic given as the sum of the numbers of left-to-right maxima and right-to-left minima exhibits the cyclic sieving phenomenon with respect to any involution with $2^{\lfloor\frac{n}{2}\rfloor}$ fixed points.
  %the Alexandersson-Kebede map.
\end{thm}

\begin{example}
  \label{ex:kappa_psi}
  Let $n = 8$. Then $\sigma = 21534687$ maps to $\psi(\sigma) = 21634587$.  Since $7$ and $8$ are already in the correct block, we switched $5$ and $6$. There is one invisible inversion in $\sigma$, $(3,5)$, and two invisible inversions in $\psi(\sigma)$, $(3, 5)$ and $(3, 6)$. In this case, $21345687$ is a fixed point.

  Next, let $n = 9$. Then $\sigma = 215346879$ maps to $\psi(\sigma) = 215346978$.  There is one invisible inversion in $\sigma$, $(3, 5)$, and two invisible inversions in $\psi(\sigma)$, $(3, 5)$ and $(7, 9)$. In this case, $215346879$ is not a fixed point, but $132456789$ is a fixed point.

\end{example}

\begin{thm}\label{thm:CSP_AK-Invisible}
  The number of invisible inversions (Statistic $1727$) exhibits the cyclic sieving phenomenon with respect to involutions with $2^{\lfloor\frac{n}{2}\rfloor}$ fixed points.
\end{thm}

\begin{proof}
  To show that the number of invisible inversions exhibits the CSP, we define a new involution $\psi: S_n \rightarrow S_n$ with $2^{\lfloor\frac{n}{2} \rfloor}$ fixed points; see Example~\ref{ex:kappa_psi}. We show the fixed points of $\psi$ have no invisible inversions, and orbits of size $2$ are made of two permutations whose number of invisible inversions differs by $1$. 

Define $\psi: S_n \rightarrow S_n$ as follows.

\begin{itemize}

    \item If $n$ is even, section the permutation into blocks of size $2$ by grouping the spots $i$ and $i + 1$ with $i$ odd. If there is some pair of values $(i, i+1)$ with $i$ odd where $i$ and $i + 1$ are not both in the $\frac{i+1}{2}$-th block, pick $i$ to be maximal and define $\psi(\sigma)$ to be the involution swapping the values $i$ and $i + 1$.

    \item If $n$ is odd, place $1$ in its own block and section the remainder of the permutation into blocks of size $2$ by grouping spots $i$ and $i + 1$ with $i$ even.  If there is some pair of values $(i, i+1)$ with $i$ even where $i$ and $i + 1$ are not both in the $\left(\frac{i}{2} + 1\right)$-th block, pick $i$ to be maximal and define $\psi(\sigma)$ by swapping the values $i$ and $i + 1$.

\end{itemize}
The fixed points are the permutations of $[n]$ where $i$ and $i+1$ are both in the $\lceil \frac{i+1}{2}\rceil$-th block for all $i$, where $i$ and $n$ have different parities. The number of such permutations is $2^{\lfloor\frac{n}{2} \rfloor}$, since ${\lfloor\frac{n}{2} \rfloor}$ is the number of blocks of size $2$.
  
  Let $\sigma$ be a permutation not fixed by $\psi$ and let $(i, i + 1)$ be the swapped pair. We show that swapping $i$ and $i + 1$ either breaks or creates an invisible inversion with those values. Consider $j, k$ such that $1 \leq j, k \leq n$, $\sigma_j = i$, $\sigma_k = i + 1$. Then $\{j , k\} \neq \{i, i+1 \}$. As $i$ is chosen to be maximal, we have $j, k \leq i+ 1$ since every value greater than $i + 1$ must already be in its assigned block to the right of spot $i + 1$. Similarly, for all $t \leq i + 1$, we have $\sigma_t \leq i + 1$.

  If $j < k$, then $\sigma_ j = i$ cannot be in its block, so $j < i$. Thus, $\psi(\sigma)_j = \sigma_k = i + 1 > \psi(\sigma)_k = \sigma_j = i > j,$ so $(j, k)$ is an invisible inversion of $\psi(\sigma)$ but not in $\sigma$ because $\sigma_j = i < \sigma_k = i + 1$. If $k < j$, then $\sigma_k = i + 1$ cannot be in its block, so $k < i$. Thus, $\sigma_k = i + 1 > \sigma_j = i > k$, so $(k , j)$ is an invisible inversion of $\sigma$ but not in $\psi(\sigma)$ since $\psi(\sigma)_k = \sigma_j = i < \psi(\sigma)_j = \sigma_k = i +1$.

  Additionally, swapping these values does not create or break any other invisible inversions. First, note that any invisible inversion not involving spots $j$ and $k$ is unaffected by this map, so we only need to consider invisible inversions involving either spot $j$ or spot $k$.  Let $\sigma_t \notin \{i, i+1\}$. Recall that if $t \leq i + 1$, then $\sigma_t \leq i + 1$, so since $\sigma_t \notin \{i, i + 1\}$, if $t \leq i + 1$, then $\sigma_t < i$.
  Hence there are no inversions of the form $(t, j)$ or $(t, k)$. Thus, there are two cases to check.

  \begin{itemize}

    \item If $(j, t)$ is an invisible inversion in $\sigma$, then $j < t$ and $\sigma_j = i > \sigma_t > j$. Then $\psi(\sigma)_j = i + 1 > \sigma_t > j$ and $(j, t)$ is an invisible inversion in $\psi(\sigma)$. Similarly, if $(j, t)$ is an invisible inversion in $\psi(\sigma)$, then $\psi(\sigma)_j = i + 1 > \psi(\sigma)_t = \sigma_t > j$ and as $\sigma_t \neq i$, $\sigma_j = i > \sigma_t > j$ and $(j, t)$ is an invisible inversion in $\sigma$.

    \item If $(k, t)$ is an invisible inversion in $\sigma$, then $k < t$ and $\sigma_k = i + 1 > \sigma_t > k$. As $\sigma_t \neq i$, then $\psi(\sigma)_k = i > \sigma_t > k$ and $(k, t)$ is an invisible inversion in $\psi(\sigma)$. Similarly, if $(k, t)$ is an invisible inversion in $\psi(\sigma)$, then $\psi(\sigma)_k = \sigma_j = i  > \psi(\sigma)_t = \sigma_t > k$, so $\sigma_k = i + 1 > \sigma_t > k$ and $(k, t)$ is an invisible inversion in $\sigma$.

\end{itemize}
Therefore, the number of invisible inversions in $\sigma$ and $\psi(\sigma)$ differs by $1$ when $\sigma$ is not a fixed point of $\psi$. Since fixed points of $\psi$ have no invisible inversions, the argument used in the proof of Theorem \ref{thm:CSP_AK} is used to prove that involutions with $2^{\lfloor \frac{n}{2}\rfloor}$ fixed points exhibit the CSP for the number of invisible inversions.
\end{proof}

\end{document}
